\chapter{Equity Index and Single-Stock Futures Worldwide}\label{ch:m1:index-and-single-stock-futures-worldwide}

In 2025 the world's exchanges traded 119 billion futures and options
contracts. Sixty-three of every hundred changed hands in Asia-Pacific,
twenty in North America, fewer than four in Europe; and eighty of every
hundred were on equities. A reader who has met only the Chicago contracts of
\cref{ch:m1:futures-contracts-and-exchanges} has seen a fifth of the
picture, measured in contracts. Measured in dollars the picture is very
different, and that gap is the first lesson of the chapter. The second is
that the equity future has relatives: contracts on the dividends an index
will pay, contracts on its total return, contracts on single shares. Each
isolates one of the terms that the fair-value formula of
\cref{prop:m1:basis-roll-and-delivery:fv} lumps together.

\section{Counting contracts, counting money}

\begin{definition}[Index point, mini and micro contracts]\label{def:m1:index-and-single-stock-futures-worldwide:point}
An \emph{index point}\index{index point} is one unit of the quoted index; a
contract's multiplier gives the point its value in money. A \emph{mini
contract}\index{mini contract} or \emph{micro contract}\index{micro contract}
is a version of an existing future with a multiplier cut, typically to a
fifth or a tenth, and then to a tenth again, so that the same market can be
traded in smaller size; minis often end up far more liquid than the
contracts they were cut from.
\end{definition}

\begin{omfigure}
\begin{tikzpicture}
\begin{axis}[width=10.5cm,height=4.8cm,xbar,bar width=9pt,
  xlabel={contracts traded in 2025 (billion)},
  symbolic y coords={Other,Europe,Latin America,North America,Asia-Pacific},ytick=data,
  tick label style={font=\small},label style={font=\small},
  xmin=0,xmax=92,xmajorgrids,grid style={gray!25},enlarge y limits=0.12,
  nodes near coords,nodes near coords style={font=\footnotesize,fill=white,inner sep=1pt,/pgf/number format/.cd,fixed,fixed zerofill,precision=1},
  table/col sep=comma]
\addplot[fill=omDef!60,draw=omDef] table[y=region,x=contracts_bn]{figdata/markets-1/22-index-and-single-stock-futures-worldwide/regions.csv};
\end{axis}
\end{tikzpicture}

\omcaption{Futures and options traded on exchanges in 2025, by region, in
number of contracts: 119.3 billion in all, 42\% fewer than in 2024. A contract
is not a unit of risk: one of them can be worth \$300\,000 and another a few
hundred dollars. Data: the industry association's annual
statistics.}\label{fig:m1:index-and-single-stock-futures-worldwide:regions}
\end{omfigure}

\begin{dated}{2026-09}{What the volume statistics say, and do not}\label{dat:m1:index-and-single-stock-futures-worldwide:fia}
The futures industry association counted 119.29 billion contracts in 2025:
88.65 billion options, down 50.3\% on the year, and 30.64 billion futures, up
8.6\%. Equity contracts were 94.79 billion, down 48.8\%, which accounts for
most of the fall. By region: Asia-Pacific 75.59 billion ($-55.6\%$), North
America 24.50 ($+24.0\%$), Latin America 11.39, Europe 4.38, others 3.43. The
counts weigh every contract equally. A halving of the world total in a year in
which futures volume rose says that the number is dominated by very small
equity options in a few markets, and by their regulators' decisions about
contract sizes and expiries, not by the amount of risk transferred.
\end{dated}

\section{Three regional benchmarks}

\begin{center}
\small
\begin{tabular}{llll}
\toprule
 & E-mini S\&P 500 & EURO STOXX 50 future & Hang Seng Index future \\
\midrule
Exchange & Chicago & Frankfurt & Hong Kong \\
Multiplier & \$50 & \euro10 & HK\$50 \\
Tick & 0.25 (\$12.50) & 1 (\euro10) & 1 (HK\$50) \\
Expiries & quarterly & quarterly & monthly and quarterly \\
Last trading day & third Friday, open & third Friday, noon & second-last trading day \\
Final settlement & opening quotation & ten-minute average & five-minute quotations, averaged \\
\bottomrule
\end{tabular}
\end{center}

\begin{dated}{2026-09}{Hours}\label{dat:m1:index-and-single-stock-futures-worldwide:hours}
The EURO STOXX 50 future trades from 02:10 to 22:00 Frankfurt time. The Hang
Seng future trades from 09:15 to 12:00 and 13:00 to 16:30 Hong Kong time, then
in an after-hours session from 17:00 to 03:00 whose trades belong to the next
day's clearing. Its contract months are the spot month, the next three
calendar months and the next three quarter months, with longer-dated June
and December contracts.
\end{dated}

Each of the three is the hedging instrument of its region's equity
derivatives and the first price to move on news when the cash market is
closed. Between them they are open almost continuously from the Asian
morning to the American evening (\cref{fig:m1:index-and-single-stock-futures-worldwide:relay}),
and a firm that trades one around the clock must know what is open
elsewhere: liquidity in each contract rises sharply when its own cash market
is open and again when another region's opens.

\begin{omfigure}
\begin{tikzpicture}
\begin{axis}[width=12cm,height=3.8cm,
  xlabel={hour (UTC), Wednesday 16 September 2026},
  ytick={1,2,3},yticklabels={ES,FESX,HSI},
  tick label style={font=\small},label style={font=\small},yticklabel style={font=\small\ttfamily},
  xmin=0,xmax=24,ymin=0.4,ymax=3.6,xtick={0,3,6,9,12,15,18,21,24},xmajorgrids,grid style={gray!25},
  unbounded coords=jump,table/col sep=comma]
\addplot[omDef,line width=7pt] table[x=hour,y=row]{figdata/markets-1/22-index-and-single-stock-futures-worldwide/relay.csv};
\end{axis}
\end{tikzpicture}

\omcaption{The relay on one day, in UTC. Hong Kong breaks for lunch; nothing
is open between 21:00 and 22:00. The European and American bars move by an
hour against UTC, and against Hong Kong, when daylight saving ends. The ES
hours are an assumption of the chapter's test data, not a quotation of the
exchange's rule. Data: the chapter's
build.}\label{fig:m1:index-and-single-stock-futures-worldwide:relay}
\end{omfigure}

\section{Dividend futures}

\begin{definition}[Dividend future]\label{def:m1:index-and-single-stock-futures-worldwide:divfut}
A \emph{dividend future}\index{dividend future} on an index settles at the
sum of the ordinary gross dividends paid by the index's constituents during
a contract period, usually a year, expressed in \omterm{def:m1:index-and-single-stock-futures-worldwide:point}{index points}. Before expiry
its price is the market's price for that stream of dividends.
\end{definition}

\begin{dated}{2026-09}{The EURO STOXX 50 index dividend future}\label{dat:m1:index-and-single-stock-futures-worldwide:fexd}
Contract value \euro100 per index dividend point; tick 0.1 point, \euro10.
Listed: the five nearest quarterly months, the next two semi-annual months
and the eight following annual December months. Trading ends at 12:00 CET on
the third Friday of the maturity month; the final settlement price is the
cumulative total of the relevant gross dividends of the index's
constituents for the contract period, as defined by the index provider.
\end{dated}

The contract exists because of a by-product. Banks that sell structured
products on an index are left long the index's future dividends for years
ahead, a risk they cannot hedge with index futures, whose dividend exposure
they already have. Selling \omterm{def:m1:index-and-single-stock-futures-worldwide:divfut}{dividend futures} transfers it. The buyers are
funds who consider the price low: a strip of \omterm{def:m1:index-and-single-stock-futures-worldwide:divfut}{dividend futures} has often
traded below any reasonable forecast of dividends for distant years, the
discount being the price the banks pay to be rid of the risk. The strip's
slope is therefore not a forecast (\cref{fig:m1:index-and-single-stock-futures-worldwide:strip}).

\begin{proposition}[Implied dividends]\label{prop:m1:index-and-single-stock-futures-worldwide:implied}
From an index future $F$ with time $T$ to expiry, spot $S$ and financing
rate $r$, the dividends the market prices between now and expiry are
\[
D^{\text{imp}} = S(1+rT) - F
\]
(reinvestment neglected). A long index future is short these dividends: if
expected dividends fall by $\delta$ points, the future's fair value rises by
$\delta$.
\end{proposition}

\begin{omfigure}
\begin{tikzpicture}
\begin{axis}[width=10cm,height=4.6cm,ybar,bar width=14pt,
  xlabel={contract year},ylabel={dividend future (points)},
  xtick=data,x tick label style={/pgf/number format/1000 sep=},
  tick label style={font=\small},label style={font=\small},
  ymin=150,ymax=178,ymajorgrids,grid style={gray!25},
  nodes near coords,nodes near coords style={font=\footnotesize,fill=white,inner sep=1pt,/pgf/number format/.cd,fixed,fixed zerofill,precision=1},
  table/col sep=comma]
\addplot[fill=omDef!60,draw=omDef] table[x=year,y=points]{figdata/markets-1/22-index-and-single-stock-futures-worldwide/strip.csv};
\end{axis}
\end{tikzpicture}

\omcaption{An illustrative strip of annual \omterm{def:m1:index-and-single-stock-futures-worldwide:divfut}{dividend futures}. The implied
growth is $+2.1\%$ for the first year and negative afterwards, $-0.5\%$ a year
over the five years: the shape that hedging supply produces, whatever
analysts forecast. The vertical axis does not start at zero. Data: the
tutorial.}\label{fig:m1:index-and-single-stock-futures-worldwide:strip}
\end{omfigure}

\section{Total return futures}

\begin{definition}[Total return future]\label{def:m1:index-and-single-stock-futures-worldwide:trf}
A \emph{total return future}\index{total return future} is a listed future
that replicates an equity \omterm{def:m1:financing:trs}{total return swap}
(\cref{def:m1:delta-one-instruments:swap}): the holder receives the index's
price return and its distributions and pays an overnight benchmark rate plus
a spread. It is quoted and traded in that spread, in basis points a year;
the exchange converts the spread into a futures price in \omterm{def:m1:index-and-single-stock-futures-worldwide:point}{index points} using
the index close, the distributions accrued and the funding accrued.
\end{definition}

\begin{dated}{2026-09}{The EURO STOXX 50 index total return future}\label{dat:m1:index-and-single-stock-futures-worldwide:tesx}
Multiplier \euro10 per \omterm{def:m1:index-and-single-stock-futures-worldwide:point}{index point}. Quoted as an annualised spread in basis
points with one decimal, minimum change 0.5 basis point. Underlyings: the
price index, the provider's distribution index, and the euro short-term
rate, flat, as funding rate. Distributions and funding accrue from launch
and are added to the price; their daily changes pass through \omterm{def:m1:clearing-and-settlement:margin}{variation
margin}. Expiries run to almost ten years: the 21 nearest quarterly months
and five further Decembers.
\end{dated}

In the ordinary index future, financing and dividends are inside one number,
the basis, and the holder is exposed to both. The \omterm{def:m1:index-and-single-stock-futures-worldwide:trf}{total return future} pays
realised dividends and charges a realised overnight rate, so that what is
left to trade is the spread alone: the term price of equity financing, the
same quantity as the richness of the roll in
\cref{def:m1:basis-roll-and-delivery:roll}, fixed for years instead of a
quarter and cleared. Banks use it to move the financing of their equity
hedges off their balance sheets; funds use it to lock in the funding cost of
a long-term position.

\section{Single-stock futures}

\begin{definition}[Single-stock future]\label{def:m1:index-and-single-stock-futures-worldwide:ssf}
A \emph{single-stock future}\index{single-stock future} is a future on one
share, settled in cash or by delivery of the shares. It is a \omterm{def:m1:delta-one-instruments:deltaone}{delta-one}
wrapper (\cref{ch:m1:delta-one-instruments}): financing inside the price, no
stock borrow needed to be short, and corporate actions handled by the
exchange's adjustment of the contract.
\end{definition}

Their fortunes are national. They trade actively in India and on the
European derivatives exchanges, where one exchange lists futures and options
on more than 1\,200 underlyings from over twenty countries, much of the
volume being dividend- and financing-driven trades between professionals
around dividend dates. In the United States the one exchange devoted to them
stopped trading in September 2020 and withdrew its registration, swaps and
options combinations having taken the business. Where a transaction tax
applies to shares and not to futures, or where shorting shares is
restricted, the \omterm{def:m1:index-and-single-stock-futures-worldwide:ssf}{single-stock future} is the natural instrument; where
neither applies, the swap wins.

\section{Tutorial: implied dividends and the strip}\label{tut:m1:index-and-single-stock-futures-worldwide:implied}

\begin{tutorial}
\textbf{Goal.} Extract implied dividends from an index future, read growth
from a dividend strip, price a \omterm{def:m1:index-and-single-stock-futures-worldwide:trf}{total return future} from its spread, and find
out which contracts are open at a given instant. \textbf{End state:} the
three data figures of this chapter.
\begin{tutsteps}
\item \textbf{Implied dividends and growth.}
\omcode{code/markets-1/22-index-and-single-stock-futures-worldwide/python/div_implied.py}{4}{11}{Dividends implied by a future, and growth implied by a strip.}\label{lst:m1:index-and-single-stock-futures-worldwide:implied}
\item \textbf{A \omterm{def:m1:index-and-single-stock-futures-worldwide:trf}{total return future} from its spread.}
\omcode{code/markets-1/22-index-and-single-stock-futures-worldwide/python/div_implied.py}{20}{24}{From the traded spread to a price in index points.}\label{lst:m1:index-and-single-stock-futures-worldwide:trf}
\item \textbf{Sessions in UTC.} Local times, time zones, sessions that cross
midnight, weeks that start on Sunday.
\omcode{code/firm/calendar/firm_calendar.py}{29}{37}{The UTC interval of a session that opens on a given local date.}\label{lst:m1:index-and-single-stock-futures-worldwide:session}
\end{tutsteps}
\textbf{What to change next.} Add each exchange's holidays for the current
year and recompute the hours of the year in which all three contracts are
open together.
\end{tutorial}

\section{Build: the global contract calendar}\label{bld:m1:index-and-single-stock-futures-worldwide:calendar}

\begin{build}
\bfield{Purpose} Strategies of the miniature firm that run across regions
must know, for any timestamp, which contracts are open, and the simulator
must refuse orders in closed ones.
\bfield{Interface} \texttt{load(path)} returning sessions;
\texttt{Session.interval(local\_date)}; \texttt{is\_open(sessions, root,
when)}; \texttt{open\_roots(sessions, when)}; \texttt{utc\_intervals(sessions,
root, day)}.
\bfield{Rules} Sessions are stored in exchange-local time with an IANA time
zone; all answers are in UTC; daylight saving comes from the zone database,
never from a hand-written offset. A session whose close is not after its
open ends on the next local day. The week is defined by the weekdays on which
a session \emph{opens}. Holidays are out of scope and documented as such.
\bfield{Data} \texttt{data/markets-1/sessions\_sample.csv}: five rows; the
ES row is an assumption and says so in the file.
\bfield{Acceptance tests} \texttt{code/firm/calendar/tests/}: the European
open moving by an hour in UTC between September and December; Hong Kong's
lunch break and its session across midnight; Friday and Sunday boundaries;
the set of open contracts at three instants; clipping to a UTC day.
\bfield{Stretch} Holidays and half-days per exchange; expiry-day early
closes from the contract master.
\end{build}

\begin{omsources}
\item FIA, \emph{ETD Volume --- December 2025} (annual totals by category and
region).
\item Eurex, product pages \emph{EURO STOXX 50 Index Futures}, \emph{EURO
STOXX 50 Index Dividend Futures (FEXD)}, \emph{EURO STOXX 50 Index Total
Return Futures (TESX)} and the circular on the transition of the TESX funding
rate to \euro STR flat; \emph{Single Stock Futures}.
\item HKEX, \emph{Hang Seng Index Futures}, contract summary.
\item US Securities and Exchange Commission, order granting OneChicago LLC's
request to withdraw from registration, \emph{Federal Register}, 18 February
2021.
\end{omsources}

\section{Exercises}

\begin{exercise}[$\star$]\label{exo:m1:index-and-single-stock-futures-worldwide:1}
Give the notional of one contract of each of the three benchmarks at index
levels of 6\,000, 5\,400 and 26\,000, in its own currency, and the value of a
1\% move.
\end{exercise}

\begin{exercise}[$\star$]\label{exo:m1:index-and-single-stock-futures-worldwide:2}
The index is at 5\,400, the three-month future at 5\,412.0 and the financing
rate 2.5\%. Give the dividends implied for the quarter.
\end{exercise}

\begin{exercise}[$\star$]\label{exo:m1:index-and-single-stock-futures-worldwide:3}
A fund buys 50 \omterm{def:m1:index-and-single-stock-futures-worldwide:divfut}{dividend futures} of December 2027 at 171.5. The year's
dividends come to 158.2 points. Give the P\&L.
\end{exercise}

\begin{exercise}[$\star\star$]\label{exo:m1:index-and-single-stock-futures-worldwide:4}
From the strip of \cref{fig:m1:index-and-single-stock-futures-worldwide:strip}
(168.0, 171.5, 169.0, 166.0, 164.5, 163.5) compute each year's implied growth
and the annualised growth over the five years. An analyst forecasts $+4\%$ a
year. Is the strip ``wrong''?
\end{exercise}

\begin{exercise}[$\star\star$]\label{exo:m1:index-and-single-stock-futures-worldwide:5}
From \cref{fig:m1:index-and-single-stock-futures-worldwide:relay}: for how
many hours of that UTC day are all three contracts open together, and during
which hour is none open? Which bars move in December, and which way?
\end{exercise}

\begin{exercise}[$\star\star$]\label{exo:m1:index-and-single-stock-futures-worldwide:6}
A fund is long 200 \omterm{def:m1:index-and-single-stock-futures-worldwide:trf}{total return futures} with two years to expiry, bought at
a spread of 40 basis points with the index at 5\,400. The spread goes to 55.
Give the P\&L and explain its sign.
\end{exercise}

\begin{exercise}[$\star\star\star$]\label{exo:m1:index-and-single-stock-futures-worldwide:7}
\emph{Coding.} With \texttt{firm\_calendar} find, for Wednesday 16 December
2026, the UTC hour in which none of the three contracts is open, and compare
with September. Then list every instant of the week of 14 September at which
the set of open contracts changes.
\end{exercise}

\begin{exercise}[$\star\star\star$]\label{exo:m1:index-and-single-stock-futures-worldwide:8}
\emph{Find the flaw.} ``Asia-Pacific trades 63\% of the world's derivatives,
so that is where 63\% of the world's risk is transferred, and where a new
trading firm should go first.'' Use
\cref{dat:m1:index-and-single-stock-futures-worldwide:fia}.
\end{exercise}

\section{Problem: The Dividend Curve}

\begin{problem}[{Weekend problem --- trading what the index will pay}]\label{pb:m1:index-and-single-stock-futures-worldwide:1}
The index is at 5\,400. Annual \omterm{def:m1:index-and-single-stock-futures-worldwide:divfut}{dividend futures} for 2026 to 2031 trade at
168.0, 171.5, 169.0, 166.0, 164.5 and 163.5 points (\euro100 a point). The
index future has a multiplier of \euro10.

\medskip\noindent\textbf{Part I --- Reading the strip.}
\begin{enumerate}
\item Give the dividend yield implied for 2027.
\item Give the implied growth from 2026 to 2027.
\item Give the annualised implied growth from 2026 to 2031.
\item Dividends have grown 4\% a year for a decade. Give the 2031 dividend on
that trend and the discount of the 2031 future to it.
\item Who is selling the 2031 future at that discount, and why?
\end{enumerate}

\medskip\noindent\textbf{Part II --- A position.}
A fund buys 200 futures of 2031 at 163.5.
\begin{enumerate}[resume]
\item Give the notional in euros.
\item Give the P\&L if 2031 dividends come in on the 4\% trend.
\item Give the P\&L if they are cut by 30\% from the 2026 level, as in a deep
recession.
\item The future is marked daily. In a crisis it falls to 110 in a month
although 2031 is five years away. Give the \omterm{def:m1:clearing-and-settlement:margin}{variation margin}, and say why a
long-term investor cannot ignore it.
\end{enumerate}

\medskip\noindent\textbf{Part III --- Dividends inside the index future.}
A desk is long 500 index futures expiring in December 2027 as a hedge.
\begin{enumerate}[resume]
\item Expected 2027 dividends are cut by 20\%. Give the change in points.
\item Give the desk's P\&L from this repricing alone, with its sign.
\item How many 2027 \omterm{def:m1:index-and-single-stock-futures-worldwide:divfut}{dividend futures} would neutralise that exposure?
\item Explain why a structured-products desk ends up long dividends.
\item Why would the same desk prefer a \omterm{def:m1:index-and-single-stock-futures-worldwide:trf}{total return future} to the ordinary
future for its hedge?
\end{enumerate}

\medskip\noindent\textbf{Part IV --- Judgement.}
\begin{enumerate}[resume]
\item Is the strip a forecast? What is it?
\item \omterm{def:m1:index-and-single-stock-futures-worldwide:divfut}{Dividend futures} fell much more than dividends were eventually cut in
past crises. What does that say about who holds them?
\item What plays the role of ``the underlying'' for the 2031 future today,
and what does that imply for hedging it?
\item Would you expect single-stock \omterm{def:m1:index-and-single-stock-futures-worldwide:divfut}{dividend futures} to be more or less
liquid than the index contract? Why do they exist?
\item State the \emph{named result}: the implied dividend growth, first year
and five-year annualised.
\item In one sentence: what does a \omterm{def:m1:index-and-single-stock-futures-worldwide:divfut}{dividend future} unbundle?
\end{enumerate}
\end{problem}

\section{Interview questions}

\begin{interviewq}[$\star$ \iqroles{trader, researcher}]\label{iq:m1:index-and-single-stock-futures-worldwide:1}
Name the main equity-index future of three regions and one way in which
their expiries differ.
\end{interviewq}

\begin{interviewq}[$\star$ \iqroles{trader, researcher, bank}]\label{iq:m1:index-and-single-stock-futures-worldwide:2}
What is a \omterm{def:m1:index-and-single-stock-futures-worldwide:divfut}{dividend future} and who uses it?
\end{interviewq}

\begin{interviewq}[$\star\star$ \iqroles{trader, bank}]\label{iq:m1:index-and-single-stock-futures-worldwide:3}
You are long an index future. A large constituent cancels its dividend. What
happens to your P\&L?
\end{interviewq}

\begin{interviewq}[$\star\star$ \iqroles{trader, bank, researcher}]\label{iq:m1:index-and-single-stock-futures-worldwide:4}
What does a \omterm{def:m1:index-and-single-stock-futures-worldwide:trf}{total return future} let you trade that an ordinary index future
does not?
\end{interviewq}

\begin{interviewq}[$\star\star$ \iqroles{developer, researcher}]\label{iq:m1:index-and-single-stock-futures-worldwide:5}
How would you store trading hours for two hundred contracts on twenty
exchanges?
\end{interviewq}

\begin{interviewq}[$\star\star\star$ \iqroles{researcher, trader}]\label{iq:m1:index-and-single-stock-futures-worldwide:6}
A league table ranks exchanges by number of contracts traded. How would you
turn it into something a trading firm can use?
\end{interviewq}
